% !TEX program = lualatex
\documentclass[aspectratio=43,professionalfonts,handout]{beamer}
\usepackage{beamer_template}

\newcommand{\LectureNum}{28}
\newcommand{\LectureTitle}{後期前半の総整理}
\newcommand{\TextPages}{p.\,102-139}
\newcommand{\NextTopic}{後期復習（三角関数）}
\newcommand{\NextPages}{p.\,140-170}
\newcommand{\CourseName}{基礎数学B}
\renewcommand{\TermName}{後期}

\begin{document}

% -----------------------------------------------
% スライド1: 表紙＋目標
% -----------------------------------------------
\begin{frame}{\CourseName\quad 第\LectureNum 回「\LectureTitle 」}
  {\small \TermName\quad \TeacherName\quad ／\quad 教科書 \TextPages\quad
  \textbf{目標}：後期前半（4章・5章1節）の内容を整理する}

\end{frame}

% -----------------------------------------------
% スライド2: 公式・性質
% -----------------------------------------------
\begin{frame}{公式・性質}
  \begin{block}{}
  \begin{itemize}
    \item 指数・対数の計算規則（$\log$ の加減と指数の積商）
    \item 正弦定理・余弦定理・面積公式の一覧
  \end{itemize}
  \end{block}
\end{frame}

% -----------------------------------------------
% まとめ
% -----------------------------------------------
\begin{frame}{まとめ}
  \begin{itemize}
    \item 指数・対数の計算規則（$\log$ の加減と指数の積商）を確認する
    \item 三角比の定理（正弦・余弦定理）と面積公式を整理する
    \item 問題の与え方に応じて正弦定理か余弦定理かを選ぶ
  \end{itemize}
\end{frame}

\end{document}
